\documentclass[10pt]{article}
\usepackage{pifont}
\usepackage[a4paper,margin=1.9cm]{geometry}
\usepackage[T1]{fontenc}
\usepackage{lmodern}
\usepackage{amsmath}
\usepackage{amssymb}
\usepackage{textcomp}
\usepackage{graphicx}
\usepackage{booktabs}
\usepackage{longtable}
\usepackage{array}
\usepackage{float}
\usepackage{xcolor}
\usepackage{caption}
\usepackage[hidelinks]{hyperref}
\usepackage[normalem]{ulem}
\newcommand{\dead}[1]{\textcolor{gray}{\sout{#1}}}
\setlength{\parindent}{0pt}
\setlength{\parskip}{3pt}
\captionsetup{font=small,labelfont=bf}
\graphicspath{{plots/}{run_figures/}{./}}


\title{\vspace{-1.4cm}\textbf{AWI--ESM3 v3.5 PI-control tuning campaign --- summary}\\[2pt]
\large Forty-five rounds, from a boreal forest that vanished to a production line that
carries one again, a Labrador Sea traced to its GM setting, and a model that is too warm for
the pre-industrial}
\author{OpenIFS 48r1 / FESOM2 / LPJ-GUESS / OASIS / XIOS, TCO95L91 / CORE3\\[1pt]
\normalsize supersedes the 2026--08--24 summary; rounds 41--46 added 2026--10--03}
\date{2026--10--03}
\begin{document}
\maketitle
\vspace{-0.7cm}

\begin{abstract}\noindent
The short reading of a campaign whose full record runs to 163 pages
(\texttt{report.tex}). Three objectives throughout: the \textbf{global radiative
imbalance}, the \textbf{Southern Ocean}, and the \textbf{NH high-latitude cold bias /
boreal forest}. \S\ref{sec:timeline} is the whole campaign in the order it happened,
rounds 1 to 46; the sections after it are the state now.

\medskip\noindent\textbf{The founding finding still stands.} Every lever that lowers the
TOA imbalance does it by cooling the model. The 29 early runs lie on a line, not in a
cloud: about $0.6$\,K of global cold per W\,m$^{-2}$ of TOA. Underneath that line sits a
structural floor that no atmospheric retune reaches --- a boreal forest that LPJ-GUESS
does not deliver, worth $-3.1$\,K of Siberian DJF on a clean with/without pair.

\medskip\noindent\textbf{The campaign has now left lever-screening behind} and produced a
single 160-year pre-industrial line, \textbf{PICAL\_ccnice} (1940--2100), on the v3.5
defaults. It is not converged, and the drift is cold: $0.214\pm0.025$\,K/decade globally
over 2060--2099, $0.511\pm0.118$ at 60--90\,N, with March NH ice extent growing.

\medskip\noindent\textbf{Positive TOA does not mean warming here.} Net TOA over 2080--2099
is $+0.167$\,W\,m$^{-2}$ while the surface cools. Total ocean heat content is flat
($+0.04$\,ZJ/decade): the upper 700\,m loses to below 700\,m at $0.32$\,W\,m$^{-2}$
equivalent, almost exactly balanced. \textbf{The cold drift is an ocean redistribution,
not an energy leak.} On one axis from 1850 to 2139 the warm layer at 300--1000\,m builds without
interruption, and of the mixing schemes it is IDEMIX that feeds 700--1500\,m ($+31$\,mK/decade
against $+9$ for TKE alone and $+11$ for KPP). A TKE-only continuation is the first run aimed at it.

\medskip\noindent\textbf{The Southern Ocean has resisted the whole campaign.} It generates
the imbalance rather than responding to it ($+0.383$ and $+0.353$\,W\,m$^{-2}$ from the two
southern bands), 84\,\% of its cloud error is already present in AMIP over a correct ocean,
and five more levers were excluded in round 37. It sits at 62\,\% of observed winter ice
volume and the ACC at 98.9\,Sv against $134\pm11.2$ observed.

\medskip\noindent\textbf{The boreal forest was gone, and the line now carries one.} BNS 0.001
and 5\,\% tree cover in the Siberian box at 2099 on ccnice, every establishment gate open.
\textbf{PICAL\_crunveg} (on albedo since 2100) restarts LPJ-GUESS from the CRUNCEP state with
the merged canopy code and keeps 81\,\% of its larch after 19 years; the forest costs about
$2\times10^3$\,km$^3$ of NH ice and $+0.24$\,K globally.

\medskip\noindent\textbf{The GM regime, a pre-industrial reference, and the Antarctic pack.} The Labrador Sea ice-capped
regime appears only under Round 36's GM 2500 / Redi 1000 (61 runs); GM 2500's heat-uptake
benefit is Southern Ocean ventilation at 34--55\,S. Of four 20-year GM arms at the production
step, \textbf{hemispheric GM} (2500 south, 1000 north) restores Labrador convection and moves
nothing else significantly; it is adopted and overshoots. \textbf{Against a pre-industrial
reference the model is $+0.46$\,K too warm}, in summer mid-latitudes, with a cold Arctic
winter: too warm and a positive imbalance now point the same way. The Antarctic winter warm ring
sits over a too-diffuse pack whose drift over-responds to winds that are about right.
\end{abstract}

\section{Scorecard against the three objectives}\label{sec:score}

\begin{center}\small
\begin{tabular}{p{2.5cm}p{4.3cm}p{4.3cm}p{4.3cm}}
\toprule
objective & at the start & what the campaign did to it & where it stands \\
\midrule
Radiative imbalance &
Coupled net TOA $+1.44$ to $+1.73$\,W\,m$^{-2}$, and every lever that lowered it cooled the
model $0.6$\,K per W\,m$^{-2}$. &
Measured that ${\approx}45\,\%$ is atmospheric (AMIP $+0.67$); re-based the target to $0$;
found that decade-4 TOA follows mid-depth heat burial ($r=+0.72$) and that GM is the lever
that reaches it. &
\textbf{Partly understood, not fixed.} gmhemi1800 2130--39: net TOA $+0.52$, surface
$+0.67$, ocean uptake $0.69$\,W\,m$^{-2}$ mostly below 700\,m, while the line still cools
$0.06$\,K/decade. Against a PI reference the model is $+0.46$\,K \emph{warm}, so the imbalance
and the temperature error now point the same way (\S\ref{sec:pibias}). \\
\addlinespace
Southern Ocean &
Largest single regional term, $+7.46$\,W\,m$^{-2}$ period-clean, and a winter over-ice warm
bias. &
Established it is cloud \emph{amount} at 60--45\,S and \emph{opacity} at 90--60\,S; that
84\,\% is inherited from AMIP; traced and priced the marine CCN path; excluded nine levers
in total. &
\textbf{Open; to be documented as structural.} 67\,\% of observed winter volume
(gmhemi1800), ACC 30\,\% weak, GM 2500 costs 16.7\,Sv. The winter warm ring is a diffuse pack,
not a missing edge, and its drift over-responds (\S\ref{sec:ice}). \\
\addlinespace
NH cold bias / boreal forest &
Siberian JJA $-2.58$\,K against ERA5; \texttt{cvh} $0.63\to0.02$ when LPJ-GUESS is switched
on. &
Split it into a universal part, an orographic part and a boreal-specific part; recovered
${\approx}1.0$\,K of the last with G4/K1 and another $0.77$ with the stable-BL lever;
closed land albedo and surface exchange as explanations. &
\textbf{Atmospheric half improved; the forest is back on the line.} crunveg keeps 81\,\% of
its larch; against a PI reference the NH summer land is now \emph{warm} ($+2.1$\,K at
30--60\,N) and the Arctic winter cold (\S\ref{sec:boreal}, \S\ref{sec:pibias}). \\
\bottomrule
\end{tabular}
\end{center}

\section{Where the model stands}\label{sec:stands}

\textbf{PICAL\_ccnice} 2090--2099 and \textbf{gmhemi1800} 2130--2139 (the crunveg line with
hemispheric GM, \S\ref{sec:gm}) unless stated. Observations as noted.

{\small
\begin{longtable}{lrrrl}
\toprule
quantity & ccnice & gmhemi1800 & observed & source \\
\midrule
\endfirsthead
\toprule
quantity & ccnice & gmhemi1800 & observed & source \\
\midrule
\endhead
net TOA (ccnice 2080--99)     & $+0.167$ & $+0.519$ & 0 for a control & --- \\
net surface (ccnice 2080--99) & $+0.357$ & $+0.673$ & 0 for a control & --- \\
T2m global trend, K/dec       & $-0.214$ & $-0.059$ & 0 for a control & 2060--99; crunveg line 2101--39 \\
T2m global, PI bias           & ---      & $+0.46$\,K & 0 & ERA5 + HadCRUT5 \\
NH April ice volume           & 34.56 & 31.69 & 31.19 & GIOMAS 1989--99 \\
NH September ice volume       & 13.97 & 12.18 & 11.96 & GIOMAS 1989--99 \\
SH September ice volume       & 12.33 & 13.33 & 19.89 & GIOMAS 1989--99 \\
SH February ice volume        & 0.73  & 1.33  & 2.07  & GIOMAS 1989--99 \\
Labrador March ice, box       & ---   & 0.10  & 0.16  & HadISST2 1979--2008 \\
Drake Passage transport       & 98.9\,Sv & --- & $134\pm11.2$ & Cunningham 2003 \\
                              &          &     & $173.3\pm10.7$ & Donohue 2016 \\
cavity basal melt             & 1062\,Gt\,yr$^{-1}$ & --- & $\sim$1100 & Adusumilli 2020 \\
                              &                     &     & $\sim$1500 & Rignot 2013 \\
CMPI (global)                 & 0.903 & 0.861 & lower is better & 2125--2139 for gmhemi1800 \\
Siberian tree cover (FPC)     & 0.050 & ---   & 0.415 & offline CRUNCEP \\
\bottomrule
\end{longtable}}

Volumes in $10^3$\,km$^3$. The NH is expected \emph{above} a satellite-era record for a
pre-industrial run; the SH is close to like-for-like because GIOMAS has no significant
Antarctic trend across it. The surface imbalance includes the snow-enthalpy term
(\S\ref{sec:traps}). gmhemi1800 is not on the same years as ccnice, so the columns differ by
land state, GM and 40 years of drift together.

\section{What v3.5 carries, and where each number came from}\label{sec:config}

The tuning now lives in \texttt{configs/setups/awiesm3/awiesm3.yaml} rather than in
runscripts, so a fresh install reproduces the line. Everything below is namelist-only
except where noted, which is what makes it mergeable with EC-Earth.

{\small
\begin{longtable}{p{4.0cm}p{2.3cm}p{9.4cm}}
\toprule
parameter & value & why, and which round decided it \\
\midrule
\endfirsthead
\toprule
parameter & value & why, and which round decided it \\
\midrule
\endhead
\texttt{RCL\_INPSEA} & 0.1 & D2b, round 11. Reduced marine ice nuclei $\to$ brighter
supercooled Southern Ocean cloud. \\
\texttt{RCL\_INPPMIN} & 70000 & Round 28 adopted 50000 (S4) and round 38 retired it: the
benefit reverses sign coupled. 70000 follows \texttt{sucldp.F90}'s own recommendation. \\
\texttt{RSBLB} & 2.0 & Rounds 31--32. The winter land lever; Arctic DJF screen $+0.614$\,K
coupled at net TOA $-0.001$. \\
\texttt{RCCNSEAICE} & 15.0 & Round 38. Open-water weighting of the whitecap sea-salt source,
which OpenIFS otherwise evaluates over sea ice. Gives the current line its name. \\
\texttt{RVRSMIN(3,4)} & 1000 & Round 14 (G4). A cap set by winter damage, not a saturation
knee. \\
\texttt{ECE\_SNOW\_SCF} & 3 (+9 coeffs) & Round 20 (P3/P5), fitted to Russian snow-course
observations after the \texttt{tanh} form was falsified against 174\,676 station soil
records. \\
\texttt{RVVEGALB} & 40 values & Round 17 (K1). Snow-free land albedo corrected against
CERES; closed land albedo as an explanation. \\
\texttt{cavity\_gamma\_scale} & 0.6 & Round 37. 1190\,Gt\,yr$^{-1}$ against 1893
uncalibrated, 1100--1325 observed. FESOM PR \#1064, merged. \\
\texttt{K\_GM\_max} / \texttt{Redi\_Kmax} & 2500 / 1000 & Round 36 (16E). Lowest TOA plateau
kept. \S\ref{sec:drake} is the bill nobody scored, and the Labrador capping another
(\S\ref{sec:gm}). \\
\texttt{K\_GM\_max\_NH} / \texttt{K\_GM\_hemi\_trans} & 1000 / 10 & Round 43 (gmhemi1800).
2500 stays in the SH, 1000 in the NH, linear over 10$^\circ$ across the equator. A regional
gate, accepted explicitly after consulting the FESOM core developers; FESOM PR \#1112 (open),
not yet in the production build. NH 1500 is the next test. \\
\texttt{mix\_scheme} & \texttt{cvmix\_TKE} \texttt{+cvmix\_IDEMIX} & Round 39, four-way
20-year comparison from one restart. \\
\texttt{use\_momix} & false & Round 38; the production line's first branch tests it. \\
\texttt{albsn} / \texttt{albsnm} & 0.80 / 0.65 & Round 38, computed from the model's own
melt seasons; FESOM ships 0.81/0.77, so these are deliberate. \texttt{albi}/\texttt{albim}
stay at FESOM's 0.70/0.68 --- darkening them was falsified, they are pond-free. \\
\bottomrule
\end{longtable}}

\textbf{Not levers at all:} \texttt{h0min}/\texttt{h0max} are hard-coded (0.5/1.5) in the
coupled \texttt{ice\_thermo\_cpl.F90}, and the namelist \texttt{h0}/\texttt{h0\_s} are read only
by the uncoupled thermodynamics, which is compiled out (round 45).

\textbf{Deliberately unset:} \texttt{RSNOWLIN2} stays at the code default 0.030. Raising it
to 0.04 is the only lever in fifty arms that moves the tropical longwave deficit
($+2.74$\,W\,m$^{-2}$, 120--127\,\% of it) and it warms the Southern Ocean too --- but it
costs $+0.42$\,W\,m$^{-2}$ globally. It remains the campaign's most interesting unadopted
parameter.

\section{The campaign in the order it happened}\label{sec:timeline}

One paragraph per phase, with the number that settled it. The full record of each is in
\texttt{report.tex}, whose \S1 carries the same order.

\subsection{The trade-off, and the floor underneath it (2026--07, rounds 1--8)}

Twenty-nine coupled runs, screened on net TOA against the ERA5 temperature bias, produced
the campaign's founding picture: \textbf{the runs lie on a line, not in a cloud}. Every
W\,m$^{-2}$ of TOA improvement cost about $0.6$\,K of global cold, from 080a at
$(1.73,+0.28)$ to 10A at $(0.57,-0.93)$; the coupling and ocean-mixing family cut the
imbalance 30--57\,\% but pushed the planet $1.7$--$2.1$\,K too cold. Centred RMSD was
$1.41$--$1.64$\,K across all 29 runs, a $\pm8\,\%$ spread against a raw-RMSD spread of
64\,\%: \textbf{the runs differ by an offset, not by a pattern}. That identified a shared
structural floor rather than a tuning deficiency, and sent the campaign looking for a
warming lever instead of another retune.

\subsection{The floor is LPJ-GUESS, on a clean pair (2026--07--14)}

A purpose-built no-LPJG reference --- same CORE3 ocean, same TCO95 atmosphere, matched
branches, HTESSEL with prescribed satellite vegetation, 30 years --- differs from the
baseline in LPJ-GUESS alone. \textbf{Siberian high-vegetation cover \texttt{cvh} goes
$0.63\to0.02$} when LPJ-GUESS is switched on (pan-boreal $0.49\to0.07$), and Siberia is
$-3.1$\,K colder in DJF. The no-LPJG run also beats the baseline on global T2m RMSD (1.514
against 1.717\,K) and on CMPI (0.738 against 0.745) at essentially unchanged imbalance, so
the warming lever is real and nearly free. The cold effect is present in every month except
high summer, which already ruled out snow-albedo masking as the whole mechanism: that is a
sunlight effect and cannot explain an equally cold polar night.

\subsection{The feedback-free AMIP reference (2026--07--27/28)}

An AMIP testbed was built so the atmosphere could be scored with the vegetation--albedo
feedback open but the coupled runaway severed. \textbf{(i) About 45\,\% of the imbalance is
atmospheric}: AMIP net TOA is $+0.67$\,W\,m$^{-2}$ against a coupled $+1.44$, and with SST
prescribed that cannot be ocean heat uptake --- the best branch at the time sat only
$0.25$\,W\,m$^{-2}$ above that floor, so the Part-I lever family was nearly exhausted.
\textbf{(ii) The boreal \emph{summer} cold bias is atmospheric too}: AMIP reproduces
essentially all of the coupled Siberian JJA bias ($-2.16$ against $-2.19$\,K) while its
\emph{annual} bias is only $-0.89$ against $-3.0$ coupled. Summer became an atmospheric
problem and September--April stayed the coupled/vegetation one.

\subsection{Rounds 10--12: a noise floor, a dead namelist, and the best boreal lever
(2026--07--30$\to$08--02)}

Three infrastructure findings changed how everything afterwards was measured. The
\textbf{4-year window was resolving nothing} in the boreal, so arms went to 46 years and a
detection threshold became mandatory. The \textbf{period offset} was measured rather than
estimated ($+0.42$\,K for Siberian JJA, against a $1.12$\,K guess that had already been used
to withdraw a conclusion). And \textbf{\texttt{NCMIPFIXYR} sat in a dead
\texttt{\&NAMECECMIP6} block}: the AMIP runs were never at 1850 forcing, GHG ran
1869$\to$1911, and Krakatoa is visible in 1883--86 in every arm. Deltas were unaffected, the
absolute baseline was offset ${\approx}0.15$\,W\,m$^{-2}$, and the coupled runs were
correctly pinned. On the science side D2b (ice nuclei) was adopted for the Southern Ocean
and F4 --- boreal surface exchange --- became the best lever to that point, with
G1 $=$ F4 $+$ D2b the best configuration.

\subsection{Where the bias is not: the vertical structure (2026--08--02)}

Before spending more levers, the column was decomposed. The Siberian cold bias is
\textbf{confined below 850\,hPa}; \textbf{the soil and skin are the least-biased parts of
the column}; and \textbf{all 20 surface types are cold in every season}, so it is not a
vegetation-tile problem. That decomposition eventually retired the surface-side lever family
as a class, and it split the boreal budget into ${\approx}0.7$\,K universal, ${\approx}0.3$
orographic and $1.3$--$1.6$ boreal-specific --- with G4 taking ${\approx}0.95$ of the last,
and the remainder not reachable by a boreal-indexed lever at all.

\subsection{Rounds 13--17: land albedo, and half of it is the snow tile
(2026--08--03$\to$05)}

Half the Siberian shortwave deficit turned out not to be cloud. Of a $+0.0154$ land albedo
excess against CERES, $+0.0074$ is the snow tile and $+0.0080$ the snow-free surface; every
forest type is inside $+0.006$ once snow is masked per cell per month. The signature is
\textbf{snow-melt timing} --- accumulation is right, the melt is a month late and 30\,\% too
slow, and May is a runaway. Round 14's G4 was adopted (Siberian JJA $+0.952$\,K, $t=7.68$,
SO SW RMSE $6.88\to4.80$) with \texttt{RVRSMIN} capped at 1000 by winter damage, and round
17 adopted K1 $=$ G4 $+$ the snow-free albedo correction ($+1.036$\,K, global T2m RMSE
$1.579\to1.524$, no Arctic energy penalty). Rounds 13 and 16 --- snow cover fraction and
skin conductivity --- were both rejected on their own predictions.

\textbf{And then land albedo was closed as an explanation.} Removing 74\,\% of a directly
measured albedo bias bought $0.089$\,K of a ${\approx}0.7$\,K target; a perfect fix buys
${\approx}0.12$\,K, one sixth. That is the shape of most of this campaign: the diagnosis was
right, the leverage was not there.

\subsection{Rounds 18--21: the snow scheme, falsified on observations and refitted
(2026--08--05$\to$07)}

The strongest summer lever built (I3, $+1.322$\,K Siberian JJA, closing the surface SW
deficit to CERES exactly) was \emph{not} adopted, because it cost $+1.72$\,W\,m$^{-2}$ at
60--90\,N. The sub-grid snow depletion scheme that HTESSEL lacks was implemented, validated
offline, and it worked --- melt-out 12 May against a satellite 11 May --- and it cost
$-16.2$\,K of coupled DJF soil for $+1.2$\,K of summer. The \texttt{tanh} functional form
was then \textbf{falsified against observations rather than preference}: its range is open
at 1, so complete cover is unrepresentable, while Russian snow courses report exactly
$10/10$ in 99.74\,\% of 14\,369 Siberian DJF surveys, and the scheme costs $-20.2$\,K
against 174\,676 station soil observations where the scheme \emph{off} is right to
$+0.8$\,K. The damage was distributional, not a matter of mean cover: DJF cover is identical
to 0.002 across the arms while the soil differs by 23\,K. P3, refitted to observations, was
adopted; P5 followed. \textbf{Direct station data (RIHMI-WDC) entered the campaign here},
and the campaign stopped depending on ERA5 --- an HTESSEL sibling --- for any soil claim.

\subsection{The forest fails by two routes (2026--08--05)}

BNE/BINE are \emph{gate-excluded}: 0.7\,\% of NE-Siberian cells clear
\texttt{tcmin\_est}$=-30$ and their LAI is exactly 0.000. BNS (larch) has \emph{no} cold
gate, clears its GDD gate on 84\,\% of cells, and is \emph{outcompeted} --- grass carries
$5.8\times$ its LAI. Since NE Siberia is ${\approx}80\,\%$ larch, \textbf{the winter-gate
mechanism does not explain the region that matters}. This distinction has been rediscovered
several times since and is why \S\ref{sec:boreal} states it explicitly.

Round 27 then measured the transfer function that the campaign's founding premise assumed:
offline 2000-year LPJ-GUESS equilibria give TREE $0.202$ / BNS $0.034$ under AMIP forcing,
$0.261/0.058$ at $+1$\,K, $0.327/0.099$ at $+2$\,K, against $0.415/0.227$ under observed
forcing. Monotonic, not saturating, and $+2$\,K recovers 59\,\% of the gap ---
\textbf{closing the cold bias should recover most of the forest}. A separate transfer test
showed \textbf{forcing alone halves boreal tree cover}: same parameters, same 10\,043
gridcells, \texttt{global.ins} byte-identical, Siberian TREEFPC $0.357\to0.163$ ($-54\,\%$)
when the spin-up forcing is swapped from CRUNCEP to what the model delivers.

\subsection{Rounds 22--23: the energy target re-based, and the Southern Ocean promoted
(2026--08--09)}

Comparing a PI run against a 2000s CERES window conflates the climate difference with the
bias. Measured period-clean, \textbf{the radiative error is $+1.23$\,W\,m$^{-2}$, about
twice what the PI framing reported}. The tropical deficit largely dissolved ($-2.51$
PI-vs-CERES against $-0.67$ period-clean); the Southern Ocean got \emph{larger}
($+7.46$\,W\,m$^{-2}$, $+0.735$ of the global mean) and became unambiguously the largest
single regional term. Its deficit is \textbf{cloud amount, 65.5\,\% of it}, zonally uniform
to $1.42$\,pp --- and no lever in the campaign moves cloud area, the best managing
$+0.63$\,pp of a $6.43$\,pp gap. The reference question was settled too: CERES and MODIS put
SO cloud area at 89.5 and 89.3\,\% against the model's 83.1, while \textbf{ERA5 reproduces
the model's bias because ERA5 is IFS} and is therefore not an independent cloud reference.
Meanwhile the global cold bias was located at the tropopause, peaking at 200\,hPa over both
poles, with the atmosphere cold \emph{and} dry ($-1.62$\,K and $-10.1\,\%$ over
700--300\,hPa) in a pair that cancels in OLR.

\subsection{Rounds 24--26: marine biogenic CCN, traced end to end and priced out
(2026--08--09$\to$10)}

The one lever that reached the Southern Ocean in the right season. Droplet number there is
\textbf{correct in the mean} ($N_d\approx73$\,cm$^{-3}$ against 60--120 observed) but
\textbf{inverted in phase}: the model is lowest in DJF because it is wind-driven, where
observations peak in DJF from biogenic sulphate, and the CRE error is worst in exactly those
months. A DMS emission climatology already shipped at TCO95 and nothing read it; it was
wired through \texttt{ocp-tool}, defaulted off, and screened. At $S=166$ it moves SO SW CRE
$-7.25$\,W\,m$^{-2}$ in DJF, 67\,\% of the deficit in the worst season, from the bloom
rather than a tuned constant. \textbf{It still cannot be adopted}: the tropical cost is
$0.0129$\,W\,m$^{-2}$ per unit $S$, so the tolerance binds at $S\approx39$, below the
published floor of 43, and scaling cannot separate the bands because tropical background CCN
is \emph{lower} than the Southern Ocean's.

\subsection{Rounds 27--28: a tuned stack nobody ran, and S4 (2026--08--10$\to$11)}

A jointly-tuned eight-parameter cloud stack for TCO95L91 was found sitting unused in the
v3.4.2 DARS2 runscripts, with no \texttt{fort.4} on disk setting any of it; it had been
invisible because \texttt{preflight} looked in one place. Taken apart, it yielded
\texttt{RSNOWLIN2}, still the only lever that moves the tropical longwave. The round also
established that \textbf{Southern Ocean levers must be scored on net, not SW CRE} ---
longwave cancels a consistent 25--29\,\% across the 44-year archive. Then S4
(\texttt{RCL\_INPPMIN} $70000\to50000$), one namelist number that had sat unevaluated for
two days, came out as the best configuration the campaign had: SO surface bias
$+3.66\to-0.96$, global net TOA $+0.64\to+0.25$. \textbf{It was the headline of the
2026--08--11 summary and is now retired} --- its benefit reverses sign coupled (Siberian JJA
$+1.194$\,K AMIP against $-0.365$\,K coupled). The habit of striking falsified claims through
rather than deleting them dates from this episode.

\subsection{Rounds 29--30: the land-surface repair, Raupach, and an inherited bias
(2026--08--13$\to$19)}

\textbf{Raupach canopy roughness is dead structurally, not statistically.} Its winter
mechanism needs a canopy, and the model's Siberian trees are $2.89$\,m tall with 2\,cm
stems, giving $z_0=0.0132$\,m against IFS's $0.058$ --- it makes the winter surface
$4.4\times$ smoother. \textbf{The lever cannot bootstrap the forest it requires.} The
land-surface repair (107 new \texttt{fail()} calls classified, peat and LUH3 reconciled to
one epoch) is a trade: $+0.308$\,K of Siberian DJF soil against a JJA bias moved from
$-0.673$ to $-1.123$. Round 30 then produced the reframing
that governs the Southern Ocean work since: \textbf{84\,\% of the coupled SO cloud error is
already present in AMIP} over a prescribed, correct ocean. It is not a coupling problem.
Round 30 also killed the ``$-1.0$\,W\,m$^{-2}$ structural leak'': it was the omitted
snow-enthalpy term, and with it the column conserves to $0.02$--$0.04$\,W\,m$^{-2}$ and a
TOA-zero run drifts $-0.02$\,K/century. \textbf{Tuning net TOA to zero does stop the ocean
drift} --- which is what makes \S\ref{sec:drift} surprising.

\subsection{Rounds 31--32: the stable boundary layer, and the campaign's best run
(2026--08--21$\to$23)}

The decomposition said the winter land bias sits between the screen and 925\,hPa: in the
best arm the air-side inversion bias is $+2.81$\,K while the skin-side bias is $+0.11$\,K.
That excluded every surface-side lever by construction and pointed at stable mixing, so
\texttt{\&NAMVDFS} was created to expose it. \texttt{RSBLB 2.0} gives Arctic DJF screen
$+0.614$/$+0.674$\,K coupled, and \textbf{it is free in energy terms coupled} --- global net
TOA $-0.001$/$-0.027$, both null, against AMIP's $+0.175$: the tropical shortwave gain is
absorbed as ocean warming and cancelled in the longwave. \textbf{Winter and summer moved
together for the first time in the campaign.} 11R scored CMPI $0.7577$ against 11E's
$0.8237$, the best in the campaign. The mixing-length route (SB3) was falsified: its
tropical SW cost per K of inversion is $3.43$ against \texttt{RSBLB}'s $0.88$, so the trade
is a property of the scheme, not of the knob. Only about half the residual land bias is
reachable by any mixing lever --- of $-2.23$\,K, ${\approx}1.0$ is excess inversion and
${\approx}1.2$ is a deep cold bias still present at 850\,hPa.

\subsection{Round 33: the Southern Ocean is two problems (2026--08--23)}

\textbf{It generates the imbalance rather than responding to it}: 90--60\,S and 60--45\,S
contribute $+0.383$ and $+0.353$\,W\,m$^{-2}$, the two largest positive terms. The two bands
need \emph{opposite} corrections --- at 90--60\,S cover is right to 2\,pp but each cloud is
8\,\% too dim, at 60--45\,S there are 5\,pp too few clouds each 9\,\% too bright --- and
every campaign lever raises opacity, which is why one band is now overcorrected while the
other is not. Cloud phase flips sign at 60\,S, and the INP family had been retired on a
measurement taken in the one band where the model is too liquid. The ice melts from above
in summer and fails to refreeze from below in autumn, with a $+0.9$\,K anomaly at exactly
100\,m that autumn mixing entrains. \textbf{The ice is not too dark} (compact-ice albedo
$0.760$--$0.778$); it is too diffuse --- no cell exceeds 90\,\% concentration in
January--March, so the area-mean albedo falls to $0.49$--$0.51$.

\subsection{Rounds 34--36: where the heat goes (2026--09--06$\to$12)}

Round 34 found four defects in the coupled ice skin (\S\ref{sec:traps}). Round 35 asked
where the heat lands: \textbf{below 300\,m}, with the top
100\,m gaining $+0.013$\,W\,m$^{-2}$ while 300--2000\,m gains $+1.388$; the tropical mixed
layer is \emph{not} too deep (43.6 against WOA18's 51.6\,m); the tropical deep warming is
\textbf{heave, not diapycnal flux}; and heat crosses density surfaces at the mid- and
high-latitude winter outcrops where the mixed layer is $+48.9$ to $+155.7$\,m too deep. The
Southern Ocean \emph{loses} heat rather than burying it, with 18.7\,\% of month-nodes mixing
to the sea floor.

Round 36 closed the loop: \textbf{decade-4 TOA follows mid-depth heat storage} across 20
runs ($r=+0.72$), \textbf{stronger GM is the lever that reaches it} (GM 1500 cuts first-decade
700--2000\,m uptake by ${\approx}30\,\%$), and 16E (\texttt{K\_GM\_max} 2500, Redi 1000)
settles at the lowest TOA plateau the campaign kept, $+0.70$\,W\,m$^{-2}$. The costs were
stated at adoption: NH March ice $+1.3$\,M\,km$^2$, global T2m $+0.21$\,K, and an AMOC that
GM 1500 strengthens and GM 2500 gives back. \textbf{Drake Passage was not among the costs
scored} (\S\ref{sec:drake}).

\subsection{Rounds 37--40: the production line (2026--09--13$\to$25)}

Round 37 spent five branches on the Southern Ocean winter bias --- salt plumes, Fox-Kemper
MLE, a deeper lead-closing floor, all three together, and a calibrated basal melt --- and
\textbf{excluded all of them}. It is a freshwater regime, not a mixing deficiency. Round 38
moved the v3.5 defaults into \texttt{awiesm3.yaml}, started the line, and computed the
sea-ice albedos from the model's own melt seasons rather than guessing: the four
\texttt{ice\_therm} albedos are \emph{pond-free}, so darkening them to stand in for ponds
double-counts, and the two snow albedos act in different hemispheres by measured surface
state ($\partial\alpha/\partial\,\texttt{albsn}$ 0.908 in SH DJF against 0.149 in NH JJA).
Round 39 scored four vertical mixing schemes and adopted \texttt{cvmix\_TKE+cvmix\_IDEMIX},
finding two real FESOM bugs on the way. Round 40 carried the line to 2100 and measured what
it had become.

\subsection{Rounds 41--46: albedo, the Labrador Sea, a pre-industrial reference, and the line as one record
(2026--09--26$\to$10--03)}

Round 41 moved to albedo and started \textbf{PICAL\_crunveg}, LPJ-GUESS restarted from the
CRUNCEP state: it keeps 81\,\% of its larch after 19 years (pre-merge runs 34 and 20\,\%), and
the forest costs about $2\times10^3$\,km$^3$ of NH ice. Round 42 chased a lost Labrador Sea.
The 1800\,s deficit in IDEMIX mixing is a real numerical effect --- the element--node--element
round trip smooths once per step, so its daily strength scales as $1/\Delta t$ --- and its fix
(FESOM PR \#1103) restored Kv but \emph{not} convection; a noise-floor analysis then showed the
build and step attributions rested on one run. A 61-run history search found the capped regime
\textbf{only under Round 36's GM 2500 / Redi 1000}. Round 43 found the GM benefit is Southern
Ocean ventilation, that neither the v3.3 ramp nor the Rossby radius can separate the Labrador Sea
from the Southern Ocean margin on CORE3, and scored four 20-year arms: \textbf{hemispheric GM
adopted}. Round 44 evaluated it (CMPI 0.861), found and fixed a factor-1000 snow-enthalpy error in
\texttt{reval}'s Gregory plot, and re-scored the temperature against a PI reference: $+0.46$\,K
\emph{warm}. Round 45 took the Antarctic pack apart: diffuse, not missing; a deficit that opens in
autumn; drift that over-responds to correct winds. McPhee ocean-to-ice flux and \texttt{h0min} were
ruled out. Round 46 put the five experiments of the production line on one axis (1850--2139, with an
AMOC series) and re-scored Round 39's four mixing arms on heat by depth: IDEMIX, not TKE, warms the
mid-depth ocean. It also took the tuning to the carbon-cycle configuration, at 6.7 times the cost.

\section{The production line, and the cold drift}\label{sec:drift}

\textbf{PICAL} (1850--1929) $\to$ \textbf{PICAL\_momixoff} (branched 1920, to 2049) $\to$
\textbf{PICAL\_ccnice} (branched 1940, to 2100). Throughput on levante, 47 nodes, 10-year
legs: about 77\,SYPD under TKE+IDEMIX, 86 under \texttt{cvmix\_TKE} alone, 75 under KPP,
with $\pm15\,\%$ decade-to-decade scatter on identical configuration.

\begin{figure}[H]\centering
\includegraphics[width=0.72\textwidth]{toa_timeseries_campaign.png}
\caption{Annual net TOA along the critical path, each run drawn to the branch point of its
successor. The two vertical marks are the mixing-scheme changes. \texttt{cvmix\_TKE} at
2060 drives TOA negative; adding IDEMIX at 2080 brings it back. The run is never far from
zero and never settles on it.}
\end{figure}

\begin{center}\small
\begin{tabular}{lrrrr}
\toprule
metric & 2060--64 & 2095--99 & trend/decade & 95\,\% CI \\
\midrule
T2m global        & 14.439 & 13.645 & $-0.214$ & 0.025 \\
T2m 60--90\,N     & $-8.143$ & $-9.842$ & $-0.511$ & 0.118 \\
T2m 60--90\,S     & $-15.650$ & $-16.875$ & $-0.238$ & 0.104 \\
NH extent, March  & 16.262 & 18.188 & $+0.522$ & 0.103 \\
net TOA           & $-0.775$ & $+0.236$ & $+0.272$ & 0.079 \\
\bottomrule
\end{tabular}
\end{center}

Round 30 showed that tuning net TOA to zero \emph{does} stop the ocean drift, so a run with
positive TOA and a cooling surface needs explaining. It resolves in the ocean: total heat
content is flat at $+0.04$\,ZJ/decade over 2060--2099, while the bands are not
($-17.8$, $-34.4$, $+19.4$, $+32.7$\,ZJ/decade for 0--100, 100--700, 700--2000 and
$>$2000\,m). \textbf{The upper 700\,m is draining into the deep ocean at
$0.32$\,W\,m$^{-2}$ equivalent.} Split by mixing era, the mechanism tracks the scheme:

\begin{center}\small
\begin{tabular}{lrrrrr}
\toprule
window & upper 0--700 & deep $>$700 & total & net TOA & net surface \\
\midrule
KPP 2040--2059            & $+0.061$ & $+0.164$ & $+0.225$ & $+0.136$ & $+0.318$ \\
cvmix\_TKE 2060--2079     & $-0.488$ & $+0.176$ & $-0.312$ & $-0.468$ & $-0.251$ \\
TKE+IDEMIX 2080--2099     & $-0.201$ & $+0.452$ & $+0.251$ & $+0.167$ & $+0.357$ \\
\bottomrule
\end{tabular}
\end{center}

All in W\,m$^{-2}$ equivalent; the surface budget includes the snow-enthalpy term, about
$0.9$\,W\,m$^{-2}$ globally, without which the budget is biased warm. \texttt{cvmix\_TKE}
alone flipped the upper ocean hard negative and the whole system lost energy. Adding IDEMIX
halved the upper-ocean loss and drove deep uptake up, which is what internal-wave-driven
diapycnal mixing should do. But the upper ocean is \emph{still} losing heat, and that is
what keeps the surface cooling. \textbf{This is the top open item.}

\textbf{The crunveg line still cools, more slowly.} Over crunveg 2101--2119 followed by ob1800
2120--2139, global T2m trends $-0.059\pm0.025$\,K/decade and 60--90\,N $-0.141\pm0.115$
($2\sigma$); with gmhemi1800 as the continuation, $-0.053\pm0.025$ and $-0.203\pm0.099$. Over
2130--2139 gmhemi1800 has net TOA $0.52$, net surface $0.67$ and ocean uptake
$0.69$\,W\,m$^{-2}$, of which $0.50$ goes below 700\,m.

\textbf{The whole line on one axis.} Stitched from PICAL, momixoff, ccnice, crunveg and gmhemi1800,
the record runs 1850--2139. The warm layer at 300--1000\,m builds through all of it; 700--1500\,m is
flat under \texttt{cvmix\_TKE} alone ($+0.3\pm1.3$\,mK/decade over 2061--2079) and warms from the leg
IDEMIX came in ($+27$, then $+28$ on crunveg and $+42$ on gmhemi1800). The AMOC at 26.5\,N spins up
to 15.8\,Sv by 1900 and holds 15--16.5\,Sv through every atmospheric and sea-ice change until the
2120 branch, where it falls to 13.9\,Sv: half of that is in ob1800 too, the rest is a northward shift
under hemispheric GM with the 20--60\,N maximum unchanged at 16.3\,Sv.

\begin{figure}[H]\centering
\includegraphics[width=\textwidth]{plots/reval_critical_path_1850-2139/Hovmoeller_temp.png}
\caption{Global-mean ocean temperature against PHC3 along the production line, 1850--2139, with the
eleven changes marked. The mid-depth warm layer builds throughout and no change along the line
reverses it; the cold layer at 50--100\,m weakens under KPP and returns with \texttt{cvmix\_TKE}.}
\end{figure}

\section{Vertical mixing: the four-way comparison}\label{sec:mixing}

Four schemes, 20 years each (2020--2039) from one restart, differing only in
\texttt{mix\_scheme}.

\begin{center}\small
\begin{tabular}{llrrr}
\toprule
arm & scheme & SH Feb vol & SH Feb ext & cavity melt \\
\midrule
PICAL\_ccnice      & KPP (control)            & ---  & ---  & 1255.6 \\
PICAL\_cvKPP       & cvmix\_KPP               & ---  & ---  & 1283.5 \\
PICAL\_cvTKE       & cvmix\_TKE               & 1.31 & 1.95 & 1174.5 \\
PICAL\_cvTKEIDEMIX & cvmix\_TKE+cvmix\_IDEMIX & 1.62 & 2.26 & 1062.1 \\
\bottomrule
\end{tabular}
\end{center}

Observed SH Feb volume about 2.07, extent about 3.0. Both KPP variants lose Antarctic
summer ice, both TKE variants grow it, and \texttt{cvmix\_KPP} reproduces the KPP control
closely enough to be the null for sea ice (not for heat, below). \textbf{Adopted}, with the cost stated at adoption: NH April
volume 38.04 rising $+2.46$/decade in the test arm, where TKE alone was flat at 35.17. On
the production line the realised rate is $+1.145\pm0.968$, about half.

\textbf{What Round 39 did not score: heat by depth.} Global-mean temperature trends of the same four
arms, 2021--2039, mK/decade:

\begin{center}\small
\begin{tabular}{lrrrr}
\toprule
arm & 0--100\,m & 100--300\,m & 700--1500\,m & 1500--2500\,m \\
\midrule
KPP (control)              & $+28$  & $+38$  & $+11.3\pm1.7$ & $-2$ \\
\texttt{cvmix\_KPP}        & $-91$  & $-52$  & $+15.0\pm0.9$ & $+3$ \\
\texttt{cvmix\_TKE}        & $-214$ & $-148$ & $+8.9\pm1.8$  & $+4$ \\
\texttt{cvmix\_TKE+IDEMIX} & $-307$ & $-128$ & $+31.4\pm1.1$ & $+9$ \\
\bottomrule
\end{tabular}
\end{center}

IDEMIX is what warms 700--1500\,m. Both TKE arms drain the upper ocean, 0.26--0.34\,K colder than
KPP in the top 300\,m after a decade, which is the heat drain of \S\ref{sec:drift} at its source.
KPP was the best of the four on this measure and the worst on Antarctic summer ice; it also kept the
heat at the surface. And \texttt{cvmix\_KPP} is 0.17\,K colder than native KPP in the top 100\,m,
unexplained. \textbf{PICAL\_crunveg\_gmhemi\_tke} (TKE without IDEMIX, 2140--2159 from gmhemi1800) is
running.

\begin{figure}[H]\centering
\includegraphics[width=0.58\textwidth]{mixing_eval_volume.png}
\caption{Hemispheric sea-ice volume across the four mixing arms, 2020--2039 from one
restart. Both KPP variants lose Antarctic summer ice where both TKE variants grow it, and
\texttt{cvmix\_KPP} tracks the KPP control closely enough to serve as the null that shows
the comparison is measuring the scheme rather than the branch.}
\end{figure}

\begin{figure}[H]\centering
\includegraphics[width=0.74\textwidth]{cavity_melt_maps_mixing.png}
\caption{Ice-shelf basal melt under the cavities, 2030--2039, per mixing scheme (top) and
the change each makes against the KPP control (bottom). Both TKE arms cool the cavities
almost everywhere around the coast, with the largest reductions on the East Antarctic
fringe and in the Amundsen/Bellingshausen sector; Ross and Filchner-Ronne interiors barely
move. Symmetric-log scale, so refreezing keeps its own sign.}
\end{figure}

\begin{figure}[H]\centering
\includegraphics[width=0.68\textwidth]{pical_run_tree_gantt.png}
\caption{The PICAL lineage and every branch taken off it, on the model calendar. The
production line is the spine; the arms hang off it at the year each was branched. The arms
are cheap and the spine is expensive, which is the intended shape --- but the spine has now
run 160 years without converging.}
\end{figure}

\section{The Southern Ocean now}\label{sec:so}

\textbf{Five more levers, all excluded (round 37).} Four 20-year branches of \textbf{PI200}
at 1560 against the PI200 control on identical years --- \textbf{spp} (salt plumes),
\textbf{mle} (Fox-Kemper), \textbf{h0} (lead-closing floor $0.5\to0.25$\,m), \textbf{all3}
--- plus \textbf{cavg06} attacking the same loop from the ice shelves. None moves the winter
over-ice bias in 20 years. The model sits in one of two freshwater regimes,
salty-and-venting or fresh-with-warm-CDW, and these levers do not cross between them. The
basal-melt calibration is worth keeping on its own terms (1-year screens gave 1038 / 1190 /
1402\,Gt\,yr$^{-1}$ at 0.5 / 0.6 / 0.7 against 1893 uncalibrated) and is upstream as FESOM
PR \#1064, but it is not a Southern Ocean fix. \textbf{SSS restoring stays excluded};
\texttt{surf\_relax\_s} is 0 and remains 0, even for Weddell convection.

\textbf{The cloud shortwave deficit remains the other half.} Across the whole coupled
campaign no lever moved it until the mixing schemes: SW CRE at 45--65\,S is $-64.44$ under
KPP, $-65.73$ under \texttt{cvmix\_TKE}, against CERES $-68.02$\,W\,m$^{-2}$.
\texttt{cvmix\_TKE} closes about a third of the gap, more than any atmospheric lever
achieved, and does it from the ocean side. That is worth recording: it is the one result
suggesting the cloud error is not purely atmospheric --- against round 30's finding that
84\,\% of it is already present in AMIP.

\section{Drake Passage, and the price of the GM tuning}\label{sec:drake}

Computed with tripyview's model-accurate transect path, which resolves the 105 mesh edges
the 66\,W line crosses and integrates the flux FESOM's own discretisation carries.

\begin{center}\small
\begin{tabular}{lrrr}
\toprule
run (1580--1599, identical restarts) & GM & resolved & \textbf{net} \\
\midrule
PI200          & resolution scaling (before) & 110.0 & \textbf{110.0} \\
PI200\_gmR1500 & Rossby cutoff, $K_{GM}^{max}$ 1500 & 103.2 & \textbf{103.2} \\
PI200\_gmR2500 & Rossby cutoff, $K_{GM}^{max}$ 2500 & 93.3 & \textbf{93.3} \\
\bottomrule
\end{tabular}
\end{center}

\textbf{16.7\,Sv, 15\,\%, monotonic in $K_{GM}^{max}$}, far outside the $\sim$1.5\,Sv
decadal scatter. The bolus term responds cleanly ($\pm1.8\to\pm2.7\to\pm3.6$\,Sv) but its
\emph{net} is zero in every arm: the eddy-induced circulation across the passage is a
closed overturning cell. On the production line the transport declines slowly, 105.1
(2040--49) to 98.9 (2090--99), with no step at either mixing-scheme switch, so the ACC
weakness is structural rather than a cvmix artefact.

GM 2500 was adopted in round 36 for the mid-depth heat uptake, with the NH ice and AMOC
costs priced. Drake Passage was not scored at the time. Both constraints are now on the
ledger and they point opposite ways.

\section{The Labrador Sea and the GM setting}\label{sec:gm}

\textbf{What capped the Labrador Sea.} Across 61 coupled runs and 468 run-decades, controlled for
global SST, the persistent ice-capped Labrador state exists only under \texttt{K\_GM\_max} 2500 /
\texttt{Redi\_Kmax} 1000 --- the Round 36 setting every PICAL run carries. GM 1000 runs stay at
0.05--0.3 March ice at every global SST from 17 to 20\,$^\circ$C; the clean pair PI200\_gmR2500 /
gmR1500 gives 0.684 against 0.280 on identical years (HadISST2 0.16). The time-step and build
attributions made first rested on one run (ob1200) and are withdrawn; the IDEMIX per-step
smoothing they pointed at is real, fixed in FESOM PR \#1103, and does not restore convection.

\textbf{Why CORE2 never did.} AWI-CM3 v3.3 switched GM off below 30\,km and its CORE2 Southern
Ocean was coarse, so the Labrador Sea got about 113\,m$^2$\,s$^{-1}$ and 55--65\,S about 2000.
CORE3 refined the Southern Ocean margin to the Labrador Sea's spacing; production GM gives 983
and 981. Neither a resolution ramp nor the Rossby radius (8.0\,km against 9.3\,km) separates them.

\textbf{Where the benefit comes from.} In three GM 1500/2500 pairs the extra GM-1500 uptake enters
through the Southern Ocean surface at 34--55\,S ($+0.19$ to $+0.26$\,W\,m$^{-2}$) and is stored in
the Indo-Pacific at 700--2000\,m: the AAIW/SAMW route. The subpolar North Atlantic releases heat.
The Labrador capping is a separate North Atlantic side effect of the same parameter.

\textbf{Four arms at 1800\,s, 2120--2139, one restart}, differences against ob1800 (GM 2500),
``$*$'' beyond the paired $2\sigma/\sqrt{N}$:

\begin{center}\small
\begin{tabular}{lrrr}
\toprule
2121--2139 & gmferr1800 & \textbf{gmhemi1800} & gmrd1800 \\
\midrule
T2m global (K)                        & $-0.447^*$ & $+0.002$ & $+0.105^*$ \\
net TOA (W\,m$^{-2}$)                 & $+0.372^*$ & $+0.032$ & $-0.034$ \\
heat 700--2000\,m (W\,m$^{-2}$)       & $+0.380^*$ & $+0.055^*$ & $-0.088^*$ \\
NH March volume ($10^3$\,km$^3$)      & $+2.459^*$ & $+0.000$ & $+0.150$ \\
SH September volume ($10^3$\,km$^3$)  & $-0.023$ & $-0.066$ & $-1.669^*$ \\
NH March extent (M\,km$^2$)           & $+0.325^*$ & $-0.657^*$ & $+0.300$ \\
SH September extent (M\,km$^2$)       & $+1.380^*$ & $+0.114$ & $-0.725^*$ \\
Labrador March MLD, 2130--39 (m)      & 1723 & 2026 & 2487 \\
Labrador March ice, 2130--39          & 0.18 & 0.10 & 0.07 \\
\bottomrule
\end{tabular}
\end{center}

ob1800 itself capped in its second decade (111\,m, ice 0.43). \textbf{gmhemi1800 is adopted}:
Labrador convection near 2050\,m with no collapses and nothing else significant except NH March
extent; it overshoots (ice 0.10--0.13 against 0.16), so NH 1500 is the next test.
\textbf{gmferr1800} (Ferreira vertical scaling) is rejected for net TOA $+0.46^*$ in its second
decade with the heat carried deep and the planet $0.5$\,K colder. \textbf{gmrd1800} (Rossby-radius
K) is rejected for its Antarctic ice loss. The hemispheric gate breaks the campaign rule that
regional differences must come from physics; it was accepted as an explicit exception after
consulting the FESOM core developers.

\section{The bias against a pre-industrial reference}\label{sec:pibias}

ERA5 1990--2014 is not a PI target. Adding the HadCRUT5 warming from 1850--1900 to 1990--2014 per
5$^\circ$ box to (model $-$ ERA5) gives a PI bias. Globally the model is $-0.36$\,K against ERA5
and HadCRUT5 warmed $0.817$\,K, so \textbf{gmhemi1800 is $+0.46$\,K too warm for the
pre-industrial}.

\begin{center}\small
\begin{tabular}{lrrr}
\toprule
PI bias (K), gmhemi1800 2125--2139 & land & ocean & \\
\midrule
JJA 30--45\,N   & $+2.21$ & $+1.42$ & NH summer mid-latitudes \\
JJA 45--60\,N   & $+2.08$ & $+1.37$ & \\
DJF 60--30\,S   & $+1.94$ & $+1.38$ & SH summer, the SO cloud deficit \\
DJF 60--75\,N   & $-0.81$ & $+0.43$ & Arctic winter cold \\
DJF 75--90\,N   & $-2.08$ & $-0.61$ & \\
ANN 30\,S--30\,N& $-0.05$ & $+0.05$ & tropics right \\
\bottomrule
\end{tabular}
\end{center}

The maps add an Antarctic winter warm ring at 60--70\,S, warm eastern-boundary upwelling coasts and
a cold blob south of Greenland. Caveats: HadCRUT5 1850--1900 coverage is thin at the poles, and the
AMIP pair showed the model under-warms Siberia between epochs, so an observation-based PI target is
not the model's own PI. \textbf{This reframes the founding trade-off}: against ERA5 the model was
too cold with a positive imbalance; against a PI reference it is too warm with a positive imbalance,
and the two point the same way.

\begin{figure}[H]\centering
\includegraphics[width=0.9\textwidth]{pi_t2m_bias_gmhemi1800_2125-2139.png}
\caption{T2m bias of gmhemi1800 2125--2139 against a PI reference, (model $-$ ERA5 1990--2014) $+$
HadCRUT5 warming 1850--1900 $\to$ 1990--2014. Stippled: no HadCRUT5 1850--1900 coverage.}
\end{figure}

\textbf{The warm NH summer land is dry and sunny.} Central Asia 40--55\,N in JJA is $+2.9$\,K against
ERA5 with 0.7 against GPCP's 1.3\,mm\,d$^{-1}$, $+20$\,W\,m$^{-2}$ surface SW against CERES, an
evaporative fraction of 0.25 and top-layer soil water 0.09; Europe is $+1.3$\,K, N America $+1.6$,
Siberia $-1.1$. LPJ-GUESS cover is used directly by OpenIFS when coupled (no double counting), and
it is sparser than the IFS default: effective JJA cover 0.31 against 0.57 over central Asia, 0.69
against 0.87 over Europe, 0.65 against 0.79 over N America. Whether the sparse vegetation causes the
dry summer or follows from it cannot be separated from this run.

\section{Sea ice: volume is the state variable}\label{sec:ice}

Against GIOMAS 1989--99 the NH sits above observations in every month, which a
pre-industrial run should, but by more than the epoch offset justifies in April; the SH is
the observed cycle scaled down by about a third in every month, and that comparison is
close to like-for-like because GIOMAS has no significant Antarctic trend across the window.
Extent, scored against OSI-SAF, is reported beside volume throughout --- with the caveat
that the record on disk has no December, 286 months rather than 312, and a December value
must not be synthesised. Extent is geometrically saturated in this configuration and has
hidden a real Arctic runaway before: across the 1950 albedo step it barely moved while volume climbed from 24 to
48$\times10^3$\,km$^3$. \textbf{Volume is the state variable that says where the pack is
going}; extent is the one that gets compared to satellites. An ice-affecting lever is
scored on volume first, and a null result in extent is not evidence that nothing happened.

\textbf{The Antarctic pack is diffuse, not missing (round 45).} In SH JJA the area where both model
and HadISST have ice is 37.7\,\% of the domain and $+1.77$\,K against ERA5, at concentration 0.74
against 0.84; ice the model misses is 2.4\,\%. The 50--90\,S area deficit opens in April--May
($-2.6$ to $-2.9$\,M\,km$^2$) and stays flat through September while the dynamic loss peaks,
consistent with round 33's autumn entrainment of the $+0.9$\,K layer at 100\,m. Against OSI-455
drift 2005--2014 the SH ice moves 1.29 times too fast; Antarctic-wide deformation is about right
(ratio 0.95) but wrong by sector, with Ross and Amundsen--Bellingshausen diverging at $+1.68$ and
$+1.26$\,\%\,d$^{-1}$ against $+0.26$ and $-0.07$. The winds are not the cause (7.2 against NCEP-2's
8.0\,m\,s$^{-1}$, the Amundsen Sea Low at the same relative depth): the ice over-responds, wind factor
2.4 against 1.4\,\%, and turns $+63^\circ$ in the Ross against $+24^\circ$ observed. Slowing it would
tighten that pack but likely cost SH extent and add Arctic volume, so it is documented as a known
bias, not tuned. A McPhee $St^*u^*$ ocean-to-ice heat flux went the wrong way in both hemispheres
(SH September $-1.25^*$, NH $+1.04^*\times10^3$\,km$^3$), and \texttt{h0min} is hard-coded in the
coupled build.

\section{The boreal forest, and what it blocks}\label{sec:boreal}

Siberian box 55--75\,N / 60--180\,E, year 2099, 684 cells:

\begin{center}\small
\begin{tabular}{lrrrr}
\toprule
state & BNS & TREEFPC & GRASSFPC & grass $>$ tree \\
\midrule
offline CRUNCEP (the restart)   & 0.227 & 0.415 & --- & 29.2\,\% \\
11E after 50 coupled years      & 0.017 & 0.083 & --- & 69.6\,\% \\
\textbf{PICAL\_ccnice 2099}     & \textbf{0.001} & \textbf{0.050} & 0.216 & \textbf{84.9\,\%} \\
\bottomrule
\end{tabular}
\end{center}

Stable between 2090 and 2099, so this is the settled state, and it is the end point of the
transient round 27 identified: the coupled model dismantles an inherited forest, 92\,\% of
the larch in 50 years. \textbf{The mechanism is competition with C3 grass, not a climate
gate}: BNS has no cold limit at all, the only temperature gate is a warm one 34\,K away,
and GDD5 clears on 91\,\% of cells in the very run where the larch dies. Questions about the
boreal forest should be answered through the grass-versus-larch competition, not through
thresholds.

This is now a practical blocker. Laszlo Hajdu's LPJ-GUESS canopy work (stem area index via
\texttt{ltos}, shrubs routed to low IFS vegetation, \texttt{iftreefracca 1}) measured
$+0.46$\,K global T2m in a run that \emph{has} a boreal forest, and improved its ERA5
comparison (bias $-0.46$ against $-0.92$\,K; RMSD 1.38 against 1.77). It acts on the tree
side of the tile. With 5\,\% tree cover here, that result is not transferable, and a
continuation would measure what the change does to a grassland.
\textbf{PICAL\_crunveg} --- identical except that LPJ-GUESS starts from the CRUNCEP
spin-up end state --- was built to answer this, \dead{and cancelled before starting when the
campaign moved to albedo}, and ran on albedo from 2100.

\textbf{Round 41: the forest survives on the merged code, and costs ice.} ccnice does not regrow
(Siberian tree FPC 0.062--0.067 per decade after 2100). crunveg keeps 81\,\% of its BNS after 19
years against 34\,\% (080a) and 20\,\% (11G) from the same CRUNCEP start, also at matched GDD5 and
grass PAR --- not yet separated from state migration or non-GDD climate. Paired against ccnice over
2101--2119, the forest gives global T2m $+0.24$\,K, 60--90\,N $+0.91$\,K and NH April / September
volume $-2.0$ / $-2.35\times10^3$\,km$^3$, putting the NH at or below GIOMAS 1989--99.

\section{Model skill: the degradation is tropical}\label{sec:cmpi}

CMPI, lower is better, across four windows of the same run:

\begin{center}\small
\begin{tabular}{lrrrr}
\toprule
 & 1980--89 & 2000--09 & 2080--89 & 2090--99 \\
\midrule
\textbf{CMPI global} & 0.844 & 0.879 & 0.902 & 0.903 \\
\midrule
nino34    & 0.966 & 1.064 & 1.124 & 1.121 \\
tropics   & 0.877 & 0.874 & 1.013 & 1.011 \\
northmid  & 0.823 & 0.834 & 0.867 & 0.869 \\
arctic    & 0.867 & 0.888 & 0.880 & 0.871 \\
antarctic & 0.718 & 0.790 & 0.717 & 0.726 \\
\bottomrule
\end{tabular}
\end{center}

The polar regions are flat ($+0.003$ arctic, $+0.008$ antarctic across 110 years). The
tropics carry the whole degradation ($+0.155$ nino34, $+0.134$ tropics). By variable the
worst movers are \texttt{tas} $+0.149$, \texttt{thetao} $+0.127$, \texttt{zg} $+0.111$.
\textbf{No mechanism is attached to this yet}, and it is not where the campaign has been
looking. For scale, the campaign's best-scoring run was 11R at 0.7577.

\textbf{On the crunveg line with hemispheric GM} (gmhemi1800 2125--2139, round 44): CMPI
\textbf{0.861} against crunveg 2101--2119 0.870 and ccnice 2101--2129 0.865, with 296 of 400
entries better than CMIP6 against 286. Gains: Ni\~no3.4 SST and T2m, NH mid-latitude sea ice,
tropical 300\,hPa wind, SH mid-latitude MLD, Antarctic ice. Losses: NH mid-latitude MLD (the
Labrador overshoot), SH mid-latitude 10--100\,m salinity, Arctic surface salinity and winds, NH
mid-latitude 1000\,m temperature.

\section{Levers, across the whole campaign}\label{sec:levers}

{\small
\begin{longtable}{p{3.4cm}p{1.5cm}p{10.9cm}}
\toprule
lever & verdict & the number \\
\midrule
\endfirsthead
\toprule
lever & verdict & the number \\
\midrule
\endhead
G4 (\texttt{RVRSMIN} cap) & adopted & Siberian JJA $+0.952$\,K, $t=7.68$; SO SW RMSE
$6.88\to4.80$. \\
K1 (snow-free albedo) & adopted & $+1.036$\,K, RMSE $1.579\to1.524$, no Arctic penalty. \\
P3/P5 (\texttt{ECE\_SNOW\_SCF} 3) & adopted & Refitted to snow courses; beats scheme-off on
bias, RMSE and JJA against 105 RIHMI stations. \\
\texttt{RSBLB} 2.0 & adopted & Arctic DJF screen $+0.614$\,K coupled, net TOA null. The
first lever to move winter and summer the same way. \\
GM 2500 / Redi 1000 & adopted & Lowest TOA plateau kept, $+0.70$\,W\,m$^{-2}$. Costs NH ice,
AMOC, and 16.7\,Sv of Drake. \\
\texttt{cvmix\_TKE+IDEMIX} & adopted, under review & Closest to the observed Antarctic summer pack;
costs NH April volume, and IDEMIX warms 700--1500\,m at $+31$\,mK/decade against $+9$ for TKE alone.
A TKE-only arm is running. \\
\texttt{cavity\_gamma\_scale} 0.6 & adopted & 1190\,Gt\,yr$^{-1}$; upstream, FESOM \#1064. \\
hemispheric GM (2500 S / 1000 N) & adopted & Labrador MLD ${\sim}2050$\,m, no collapses; T2m,
TOA, Siberia, ice volume null; NH March extent $-0.66^*$. Overshoots: ice 0.10--0.13 vs 0.16.
FESOM \#1112 (open). \\
CRUNCEP land state (crunveg) & adopted & BNS 81\,\% retained after 19 years; $-2\times10^3$\,km$^3$
NH ice, $+0.24$\,K. \\
\midrule
Ferreira GM scaling & rejected & Most realistic Labrador ice (0.17--0.18), but TOA $+0.46^*$ in
decade 2, heat carried deep, global $-0.5$\,K. \\
Rossby-radius GM & rejected & SH September volume $-1.67^*\times10^3$\,km$^3$, extent $-0.73^*$. \\
GM 1500; v3.3 GM ramp & not adopted & Both open the Labrador Sea and give back the uptake benefit
(TOA $+0.38^*$, $+0.44^*$); the ramp loses 40\,\% of SH September volume. \\
IDEMIX smoothing \texttt{dtref} & fix, not a lever & Removes ${\sim}85\,\%$ of the
1800\,s mixing deficit; Labrador 414 vs 370\,m, no convection. FESOM \#1103 (open). \\
McPhee $St^*u^*$ ice--ocean flux & rejected & SH September $-1.25^*$, NH $+1.04^*\times10^3$\,km$^3$:
wrong way in both hemispheres. \\
\texttt{h0min}, \texttt{h0}/\texttt{h0\_s} & not a lever & Hard-coded in the coupled build; the
namelist values are inert. \\
slowing the Antarctic ice drift & not tuned & Winds right, ice over-responds; would likely cost SH
extent and add Arctic volume. \\
\midrule
S4 \texttt{RCL\_INPPMIN} 50000 & \textbf{retired} & Headline of the 2026--08--11 summary.
Siberian JJA $+1.194$\,K AMIP, $-0.365$\,K coupled. \\
I3 & rejected & Strongest summer lever built ($+1.322$\,K) at $+1.72$\,W\,m$^{-2}$ at
60--90\,N. \\
\texttt{tanh} snow depletion & falsified & $-20.2$\,K against 174\,676 station soil
observations; unrepresentable complete cover. \\
\texttt{LAMSK}, moss proxy, D1 & falsified & Skin conductivity $+0.125$\,K against a
$0.41$ threshold; the moss proxy moved monotonically the wrong way; \texttt{RCAPDCYCL} 4
null. \\
Raupach $z_0$ & dead structurally & Needs a canopy; makes the winter surface $4.4\times$
smoother as the model stands. \\
DMS / marine CCN & priced out & Best SO lever found ($-7.25$\,W\,m$^{-2}$ DJF) but the
tropical tolerance binds below the published coefficient floor. Reconsidered as a global cooling
lever in round 44 and declined for the last round. \\
\texttt{RSNOWLIN2} 0.04 & not adopted & Only lever that moves the tropical longwave
($+2.74$); costs $+0.42$ globally. \\
SB3 mixing length & falsified & $3.43$ tropical SW per K against \texttt{RSBLB}'s $0.88$. \\
salt plumes, MLE, \texttt{h0min} $\downarrow$, all three & excluded & Round 37, no SH winter
response in 20 years. \\
\texttt{h0min} $0.5\to1.0$ & reverted & Flux response as designed (1.71$\to$1.10), state did
not follow; 60--90\,S bias $+0.82\to+1.45$\,K. \\
darkening \texttt{albpnd} & falsified & Wrong direction, and double-counts the ponds. \\
SSS restoring & excluded by policy & Permanently. \\
\bottomrule
\end{longtable}}

\section{Reversals on the record}\label{sec:reversals}

Kept because a summary that only lists what survived misrepresents how the campaign
learned.

\begin{itemize}\itemsep1pt
\item \dead{S4 is the best configuration the campaign has.} Its benefit reverses sign
coupled. The 2026--08--11 summary led with it.
\item \dead{The period offset is ${\approx}1.12$\,K, so the boreal bias is
${\approx}-1.0$\,K.} Measured directly at $+0.42$\,K; the bias is ${\approx}-2.0$. The
method error was applying an \emph{observed} epoch change to a model that does not reproduce
it.
\item \dead{There is a $-1.0$\,W\,m$^{-2}$ structural leak in the coupled column.} It was
the omitted snow-enthalpy term.
\item \dead{The boreal problem is competition, not climate.} Both, and the competition is
driven by the forcing: warming recovers 59\,\% of the gap at $+2$\,K while
\texttt{grass$>$tree} stays at 43--45\,\%.
\item \dead{Cloud amount is unreachable, so the SO needs an opacity lever.} Round 33 showed
opacity is the right correction for 90--60\,S and the \emph{wrong} one for 60--45\,S, which
every campaign lever had been applying.
\item \dead{The ice is too dark.} It is too diffuse: compact-ice albedo is
$0.760$--$0.778$, but no cell exceeds 90\,\% concentration in January--March.
\item \dead{Positive net TOA means the run is warming.} It cools at 0.21\,K/decade.
\item \dead{The 2120s Labrador collapse is the new build and the 1800\,s IDEMIX step.} Every such
difference was against one run (ob1200); the capped regime is GM 2500 / Redi 1000.
\item \dead{The GM 2500 uptake benefit survives switching GM off on the fine mesh.} The v3.3 ramp
gave it all back and lost 40\,\% of SH September volume.
\item \dead{The model is too cold.} Against ERA5, yes ($-0.36$\,K); against a PI reference it is
$+0.46$\,K warm.
\item \dead{OpenIFS double-counts LPJ-GUESS vegetation cover.} \texttt{RVCOV} is applied only when
LPJ-GUESS is not coupled.
\item \dead{The steppe LAI of 4.9 is too high.} Per-tile LAI; cover-weighted 0.82.
\item \dead{The Antarctic winter warm ring is a missing ice edge.} It is a diffuse pack; the
missing-ice class is 2.4\,\% of the area.
\item \dead{The surface imbalance is 1.51\,W\,m$^{-2}$.} \texttt{reval} dropped the snow enthalpy;
corrected 0.67.
\item \dead{\texttt{cvmix\_KPP} is the null of the mixing comparison.} For sea ice; its top 100\,m
is 0.17\,K colder than native KPP.
\end{itemize}

\section{Bugs and traps that cost this campaign}\label{sec:traps}

\begin{center}\small
\begin{tabular}{p{4.2cm}p{11.8cm}}
\toprule
what & the lesson \\
\midrule
A dead namelist block &
\texttt{NCMIPFIXYR} sat in \texttt{\&NAMECECMIP6}, which the runtime does not read, so the
AMIP runs were never at 1850 forcing. Krakatoa in 1883--86 is the fingerprint. \\
\addlinespace
Two OASIS traps &
\texttt{oasis\_put} returns the waitgroup code for every field in a group but the last, so
the ice-skin anchor was 0.0 in all 63\,923 logged events and the skin reached
$-36\,470$\,K. And \texttt{GSMART} on a sign-changing field gave 9105 excursions above
1\,kW\,m$^{-2}$ onto real ice in one coupled year, where \texttt{GSSMAR} and
\texttt{GLOBAL} give none. \\
\addlinespace
Implicit netCDF fill value &
\texttt{read\_other\_NetCDF} left \texttt{miss} at 0.0 when a variable declared no
\texttt{missing\_value}, masking 146\,694 real zeros and admitting
$9.96921\times10^{36}$. FESOM PR \#1087. \\
\addlinespace
\texttt{MPI\_DOUBLE\_PRECISION} on a \texttt{real(WP)} buffer &
Five \texttt{MPI\_AllREDUCE} calls read 8 bytes from a 4-byte variable in a
single-precision build. Not only a diagnostic: the corrupt total was fed back into the
physics through \texttt{idemix\_botforc\_Etot}. Same PR. \\
\addlinespace
The staged library is a symlink &
\texttt{<exp>/bin/fesom/lib/fesom/libfesom.so} points into esm\_tools' per-leg snapshots
and \texttt{md5sum} follows it. The \texttt{fesom} executable is a driver stub,
byte-identical across builds, so a matching binary checksum proves nothing. Use
\texttt{ls -l} before \texttt{md5sum}. Cost a 47-node slot and a NaN'd ocean. \\
\addlinespace
The instruction file nothing reads &
\texttt{guess.ins} is hard-wired in \texttt{framework.cpp:1063} and imports
\texttt{ecearth.ins}; \texttt{global.ins} and \texttt{run\_coupled\_4\_1\_2.ins} are staged
and never reached. The two had drifted apart, \texttt{iftreefracca} 1 against 0. A
parameter put in the wrong one is silently inert. \\
\addlinespace
IDEMIX smooths once per step &
The element--node--element round trip of the internal-wave energy is an un-$\Delta t$-scaled
smoother, so 1800\,s mixes less than 1200\,s; the tuned state leans on it (Kv $-12$ to $-19\,\%$
without it). FESOM PR \#1103. \\
\addlinespace
$L_f$ in the wrong units &
\texttt{reval}'s Gregory plot multiplied snowfall in metres by 334000 instead of
$\rho_wL_f=3.3355\times10^8$, dropping a 0.81--0.86\,W\,m$^{-2}$ term and reading 1.51 for 0.67.
\texttt{release\_evaluation\_tool2} PR \#7. The same term behind round 30's ``structural leak''. \\
\addlinespace
Inert lead-closing namelist &
\texttt{h0}/\texttt{h0\_s} are read only by \texttt{ice\_thermo\_oce.F90}, compiled out when
coupled; \texttt{ice\_thermo\_cpl.F90} hard-codes \texttt{h0min} 0.5. And the
\texttt{/ice\_therm/} read (\texttt{ice\_init.F90:123}) never checks \texttt{iostat}, so one
unknown name silently resets every later name in the group. \\
\addlinespace
One realisation &
A single 1200\,s arm sitting at the 28th percentile of the null made the build and step look
guilty for two days. Test against the null before attributing. \\
\addlinespace
Chained legs and the date file &
Legs read a per-leg archive, not the source runscript, so edits need \texttt{-U}; and the
date file advances past a crashed leg, so it is not evidence of progress. \\
\addlinespace
REcoM brings its own namelists &
\texttt{bin\_2p3z2d} replaces \texttt{namelist.tra} and \texttt{namelist.io}: surface salinity
restoring comes back on and 37 production outputs disappear. And the CMIP7 XIOS templates define no
BGC field, so a carbon-cycle leg with \texttt{cmip7\_cmor\_output} writes no REcoM diagnostics. Diff
the resolved namelists against the physics run before submitting. \\
\bottomrule
\end{tabular}
\end{center}

\section{Closed by measurement}\label{sec:closed}

Claims that were live at some point and are now settled, each with the observation that
closed it.

{\small
\begin{longtable}{p{6.6cm}p{9.4cm}}
\toprule
claim & what closed it \\
\midrule
\endfirsthead
\toprule
claim & what closed it \\
\midrule
\endhead
The cold bias is a tuning artefact &
Centred RMSD is $1.41$--$1.64$\,K across all 29 early runs, a $\pm8\,\%$ spread. The runs
differ by an offset. \\
Land surface albedo explains the universal bias &
Removing 74\,\% of a measured bias bought $0.089$\,K of ${\approx}0.7$\,K. \\
Boreal needleleaf is excluded by a cold gate &
True for BNE (0.7\,\% of cells clear it) and false for BNS, which is ${\approx}80\,\%$ of
NE Siberia and is outcompeted instead. \\
ERA5 is an adequate cloud reference for the Southern Ocean &
ERA5 is IFS and reproduces the model's own bias; CERES and MODIS agree with each other
against both. \\
The tropical deep warming is diapycnal flux &
It is heave: $0.42$ of the $+0.477$\,K at fixed depth is the isopycnal sinking $+42.4$\,m. \\
The SH winter bias is a vertical-mixing deficiency &
Four levers, four nulls, 20 years each from one restart. \\
The ice is too dark because the ponds are too bright &
The \texttt{ice\_therm} albedos are pond-free; darkening them double-counts. \\
The albedo sensitivities are fixed properties &
NH $\partial\alpha/\partial\,\texttt{albsn}$ moved 0.149 $\to$ 0.504 as the Arctic cooled. \\
Positive net TOA means the run is warming &
It cools at 0.21\,K/decade; the ocean redistributes rather than accumulates. \\
GM tuning is free at the Southern Ocean &
16.7\,Sv of Drake Passage transport. \\
GM tuning is free in the North Atlantic &
The capped Labrador regime exists only under GM 2500 / Redi 1000 (61 runs; clean pair 0.684
against 0.280). \\
The Labrador collapse is an IDEMIX time-step effect &
The smoothing fix removed 85\,\% of the Kv deficit; convection stayed at 414 against 370\,m. \\
The GM uptake benefit is in the North Atlantic &
Three pairs: it enters at 34--55\,S and is stored in the Indo-Pacific at 700--2000\,m. \\
The Rossby radius separates the Labrador Sea from the SO &
8.0 against 9.3\,km; gmrd1800 lost $1.67\times10^3$\,km$^3$ of SH ice. \\
The model is too cold for a PI control &
$+0.46$\,K against an epoch-corrected reference; only the Arctic winter is cold. \\
The Antarctic drift error comes from the winds &
Wind speed and the Amundsen Sea Low are close to NCEP-2; the ice over-responds (2.4 against
1.4\,\%). \\
\bottomrule
\end{longtable}}

\section{Open, in priority order}\label{sec:open}

Re-ranked 2026--10--03 after rounds 41--46.

\begin{enumerate}\itemsep2pt
\item \textbf{The Labrador overshoot under hemispheric GM}: March ice 0.10--0.13 against 0.16, and
NH mid-latitude MLD the largest CMPI loss. NH \texttt{K\_GM\_max} 1500 waits behind the TKE-only run.
\item \textbf{The line still cools} ($-0.059\pm0.025$\,K/decade over 2101--2139) with
0.69\,W\,m$^{-2}$ going into the ocean, mostly below 700\,m. The upper-ocean drain of rounds 39--40
is the mechanism, and IDEMIX carries the mid-depth half of it (\S\ref{sec:mixing}).
\textbf{PICAL\_crunveg\_gmhemi\_tke} is the first run aimed at it; on the old line TKE alone lost
energy overall (net TOA $-0.47$\,W\,m$^{-2}$).
\item \textbf{The warm, dry NH summer land}: central Asia $+2.9$\,K, half the observed rain,
LPJ-GUESS cover 0.31 against IFS's 0.57. Cause and effect not separable from one run.
\item \textbf{The Antarctic at 67\,\% of observed winter volume}, a diffuse pack whose deficit opens
in autumn and whose drift over-responds; documented as structural.
\item \textbf{The Southern Ocean cloud amount deficit}, 84\,\% of it inherited from AMIP, now also
the $+1.4$\,K SH summer ocean PI bias. Structural.
\item \textbf{The ACC 30\,\% weak}, GM 2500 in the SH keeps the 16.7\,Sv cost. Unreconciled.
\item \textbf{The stratocumulus coasts} are warm in the PI bias map. Structural.
\item \textbf{Clean attribution of the forest's survival}: an offline old-against-new LPJ-GUESS
binary pair on identical forcing.
\end{enumerate}

\dead{Does the canopy work do anything without trees? \texttt{PICAL\_crunveg} was built to answer
it and cancelled before starting.} It ran; the forest survives (round 41).
\dead{Will the CRUNCEP state load on \texttt{TCO95-land}?} It does: crunveg ran from it.

\section{What would change the picture}\label{sec:falsifiers}

Stated in advance, so that a result which contradicts them is recognised as such.

\begin{itemize}\itemsep2pt
\item \textbf{If the upper-ocean heat drain stops without the surface warming}, then the
cold drift is not driven by the redistribution and the diagnosis in \S\ref{sec:drift} is
wrong.
\item \textbf{If NH April volume flattens on its own} over the next 50 years, the cost
priced against TKE+IDEMIX at adoption was an artefact of the 20-year test arm and the
scheme is cheaper than believed.
\item \textbf{If a Southern Ocean lever moves the winter bias in under 20 years}, then the
two-regime picture is wrong, or the regime boundary is closer than the five excluded
levers suggest.
\item \textbf{If Drake Passage recovers without GM being reduced}, the 16.7\,Sv attribution
is confounded with something else that changed between \texttt{PI200} and the arms.
\item \textbf{If the canopy work produces $+0.46$\,K here}, with 5\,\% tree cover, then it
is not acting through the tree canopy and the mechanism is misattributed.
\item \textbf{If the tropical CMPI degradation appears in a run without the cold drift},
the two are independent and the tropics need their own investigation.
\item \textbf{If closing the remaining cold bias does not recover the forest}, the round-27
transfer function does not transfer coupled, and the campaign's founding premise fails at
the last step.
\item \textbf{If a second realisation of a GM 1000 NH arm caps the Labrador Sea}, or GM 2500 runs
convect in a second realisation, the GM attribution of round 42 rests on noise as the build
attribution did.
\item \textbf{If the crunveg forest decays like 080a/11G on an offline old-binary rerun with the
same forcing}, its survival is the climate, not the merged canopy code.
\end{itemize}

\section{The run inventory since the last summary}\label{sec:runs}

Every coupled arm that produced a scoreable result since 2026--08--24; the 2120 branches all
start from crunveg's 2120 restart. All are 20 years
unless stated, all branch from one restart with the control on identical years. The
campaign's earlier arms --- the 29 screening runs, the AMIP A/D/F/G/I/K/L/P/S series and
the 11-series coupled roster --- are catalogued in \texttt{report.tex}, appendix
\texttt{app:params}.

{\small
\begin{longtable}{p{3.3cm}p{2.3cm}p{10.2cm}}
\toprule
run & years & the one thing it changed, and what it showed \\
\midrule
\endfirsthead
\toprule
run & years & the one thing it changed, and what it showed \\
\midrule
\endhead
\textbf{PI200\_spp / \_mle / \_h0 / \_all3} & 1560--79 & Salt plumes, Fox-Kemper MLE,
lead closing \texttt{h0min} $0.5\to0.25$\,m, and all three together. No SH winter response
in any, and no combination effect. \\
\textbf{PI200\_cav05/06/07} & 1600 & One-year basal-melt screens: 1038 / 1190 /
1402\,Gt\,yr$^{-1}$ against 1893 at 1.0. 0.6 taken forward. \\
\textbf{PI200\_cavg06}   & 1600--19 & \texttt{cavity\_gamma\_scale} 0.6. Melt calibrated;
no Southern Ocean response. \\
\textbf{PI200\_gmR1500}  & 1580--99 & Rossby cutoff, $K_{GM}^{max}$ 1500. Drake 103.2\,Sv. \\
\textbf{PI200\_gmR2500}  & 1580--99 & Rossby cutoff, $K_{GM}^{max}$ 2500. Drake 93.3\,Sv;
the adopted setting, and the 16.7\,Sv cost. \\
\addlinespace
\textbf{PICAL\_v35def} & 2020--49 & \texttt{RCL\_INPPMIN} 50000 vs 70000, paired with
momix-off. Context for the TOA figure. \\
\textbf{PICAL\_v35def\_darkice} & 2050--99 & \texttt{albi}/\texttt{albim} darkened against
its parent. An Arctic runaway fifty years did not arrest. \\
\addlinespace
\textbf{PICAL\_cvKPP}    & 2020--39 & \texttt{cvmix\_KPP}. The useful null: reproduces the
KPP control. \\
\textbf{PICAL\_cvTKE}    & 2020--39 & \texttt{cvmix\_TKE}. Grows the Antarctic summer pack;
costs energy. \\
\textbf{PICAL\_cvTKEIDEMIX} & 2020--39 & \texttt{cvmix\_TKE+cvmix\_IDEMIX}. Closest to the
observed summer pack; adopted. \\
\addlinespace
\textbf{PICAL\_ccnice}   & 1940--2129 & The production line to 2099; continued on albedo to
2129 as the unforested control. \\
\textbf{PICAL\_crunveg}  & 2100--19 & LPJ-GUESS from the CRUNCEP state with the merged canopy code.
BNS 81\,\% retained; the new production line. \\
\addlinespace
\textbf{gm2500} (\texttt{\_ob1200}), \textbf{ob1800} & 2120--39 & Old build at 1200 / 1800\,s
(ob1200 to 2129). The GM 2500 references. \\
\textbf{ts1200}          & 2120--29 & New build at 1200\,s. Labrador 520\,m. \\
\textbf{ihf0 / ihf1}     & 2120--29 & New build at 1800\,s, McPhee flux off/on. Wrong way in both
hemispheres. \\
\textbf{nx1800}          & 2120--29 & ihf0 + IDEMIX smoothing fix. Kv restored, Labrador 414\,m;
second leg cancelled. \\
\textbf{gm1500 / gmramp} & 2120--29 & 1200\,s, GM 1500 / v3.3 ramp. Both open the Labrador Sea and
give back the uptake benefit; second decades cancelled. \\
\textbf{gmferr1800}      & 2120--39 & Ferreira scaling. Rejected, TOA $+0.46^*$. \\
\textbf{gmhemi1800}      & 2120--39 & Hemispheric GM. Adopted. \\
\textbf{gmrd1800}        & 2120--39 & Rossby-radius GM. Rejected, SH ice. \\
\textbf{PICAL\_cc}       & 2140--44 & gmhemi1800 tuning on \texttt{develop-cc}: REcoM cold start,
concentration-driven PI CO$_2$. 0.674\,s per model hour; no BGC diagnostics in this leg; paused. \\
\textbf{gmhemi\_tke}     & 2140--59 & gmhemi1800 with \texttt{cvmix\_TKE}, IDEMIX off. Running. \\
\bottomrule
\end{longtable}}

\section{The evaluation protocol, and which references are independent}\label{sec:eval}

Every window is scored the same way, so results from different rounds are comparable.

{\small
\begin{longtable}{p{3.2cm}p{12.6cm}}
\toprule
target & reference, and its caveat \\
\midrule
\endfirsthead
\toprule
target & reference, and its caveat \\
\midrule
\endhead
Sea-ice volume & \textbf{GIOMAS}. A model with assimilation, not a measurement: Arctic
volume carries 10--20\,\% uncertainty and Antarctic more. \textbf{Use 1989--99} for a PI
run --- the Arctic thinned sharply across the full record, the Antarctic has no trend. \\
Sea-ice extent & \textbf{OSI-SAF}. The record on disk has \textbf{no December} --- 286
months --- and a December value must not be synthesised. \\
Radiation & \textbf{CERES}. Period offsets matter, and measuring one rather than estimating
it once reversed a conclusion. \\
Screen temperature & \textbf{ERA5}, except for anything the IFS land surface decides: ERA5
is an HTESSEL sibling and reproduces the model's own cloud and snow biases. \\
Soil temperature & \textbf{RIHMI-WDC} station data, 174\,676 observations. This is why the
campaign no longer depends on ERA5 for a soil claim. \\
Snow cover & \textbf{Rutgers} and Russian snow courses (14\,369 Siberian DJF surveys). \\
Cloud amount & \textbf{CERES} and \textbf{MODIS}, which agree with each other, over dark
ocean only --- over ice they disagree by 7.7\,pp. \\
Ocean T/S & \textbf{PHC3} on the \emph{new} CORE3 mesh (220509 nodes). The 211567-node beta
mesh climatology is a different file and mixing them fails silently. \\
Drake Passage & \textbf{Cunningham 2003} $134\pm11.2$\,Sv and \textbf{Donohue 2016}
$173.3\pm10.7$\,Sv. Donohue includes near-bottom flow a 47-level mesh cannot carry, so the
gap to 134 is the real bias. \\
Cavity basal melt & \textbf{Rignot 2013} $\sim$1500 and \textbf{Adusumilli 2020}
$\sim$1100\,Gt\,yr$^{-1}$; the calibration target is 1100--1325. \\
Composite skill & \textbf{cmpitool} against ERA5, thirteen variables by region and season.
Useful for regressions, blunt for mechanisms. \\
PI temperature & \textbf{ERA5 1990--2014 $+$ HadCRUT5} warming 1850--1900 $\to$ 1990--2014 per
5$^\circ$ box (0.817\,K global). Coverage is thin at the poles, and the model's own epoch warming
differs (it under-warms Siberia). \\
Sea-ice drift & \textbf{OSI-455} v1.0 daily drift 2005--2014, satellite vectors only, with
\textbf{NCEP-2} 10\,m winds. A positive-SAM epoch, and satellite drift misses fast MIZ ice. \\
Labrador Sea ice & \textbf{HadISST2} March 1979--2008: 0.16 in the 56--62\,N 60--50\,W box, 0.37
over the basin; its pre-1979 values are a flat climatology. \\
\bottomrule
\end{longtable}}

\textbf{How long an arm has to run.} The standing threshold is $2\sigma\sqrt{2/N}$ from the
control's own scatter in the same window, with no autocorrelation correction, so it is a
floor rather than a true confidence interval --- and branched arms are scored
\emph{paired}, because they share a restart and the unpaired form discards that and once
hid a resolved result. Measured on this configuration: Drake Passage scatter $\sim$1.5\,Sv;
ice volume needs 20 years to resolve $\pm1\times10^3$\,km$^3$/decade; decadal TOA means are
stable to $\sim\pm0.1$\,W\,m$^{-2}$; wallclock varies $\pm15\,\%$, so SYPD comparisons need
more than one leg each. The per-band one-year thresholds from round 26 remain in use:
SO SW CRE annual $\pm1.97$, SO DJF $\pm3.82$, tropics annual $\pm0.99$, global $\pm1.11$.

The full per-window diagnostic set is \texttt{reval}
(\texttt{release\_evaluation\_tool2}), driven by one config per window. Copying
\texttt{configs/awiesm3\_v35\_pical\_ccnice\_2090s.py} and changing four fields
(\texttt{model\_version}, \texttt{pi\_ctrl\_start}/\texttt{end}, \texttt{spinup\_end},
\texttt{out\_path}) is the whole procedure; it submits 26 jobs and an HTML report.

\section{Method notes}\label{sec:method}

Techniques that earned their keep, recorded because they generalise.

\textbf{Measure the reference-period offset, do not estimate it.} Applying an
\emph{observed} epoch change to a model that does not reproduce that change put the boreal
bias off by a factor of two and cost a withdrawn conclusion. The AMIP present-day run
measures it directly.

\textbf{Verify the artefact that executes, not the one that was edited.} A namelist
parameter is live only if the \emph{staged} shared library carries it, and the file that is
read may not be the one that looks live --- LPJ-GUESS reads a hard-wired \texttt{guess.ins}
while several plausible instruction files are staged beside it and never opened. Both traps
cost a run each.

\textbf{Prove a physical diagnostic against the input file's own metadata.} The IDEMIX gate
works because $0.9179869$\,TW is the forcing file's \texttt{Etot\_in\_TW\_bin} attribute,
not a remembered number. A self-consistency check would have passed on the broken build.

\textbf{Use the model's own integrals where they exist.} Hemispheric ice volume comes from
FESOM's \texttt{sivoln}/\texttt{sivols}, computed on the native mesh inside the model, so it
cannot pick up the wrong mesh. Likewise Drake Passage is computed on the model's own transect
path rather than by interpolating to a straight line, which would not conserve.

\textbf{Screen at one year, or offline, before spending forty.} The basal-melt calibration
was settled by three one-year runs; the snow scheme was validated offline before
implementation; the forcing-transfer test used two spin-ups that already existed.

\textbf{Keep one store, append to it.} \texttt{data/coupled\_annual\_diag.nc} holds annual
TOA, surface and four-band ocean heat content for 90 runs across 750 model years, and
survives the deletion of the raw output it was built from. Several results here would be
unrecoverable otherwise.

\section{The software stack, and the move to albedo}\label{sec:stack}

The production runs of rounds 41--45 used the old FESOM build plus local branches (GM arms) or the
new build (\texttt{ts1200}, \texttt{ihf*}, \texttt{nx1800}). For the last round a fresh
\texttt{awiesm3-develop-cc} was installed on 2026--10--02/03 at
\path{/albedo/work/user/jstreffi/model_codes/awiesm3-develop-cc} with \texttt{esm\_master
comp-awiesm3-develop-cc}:

\begin{center}\small
\begin{tabular}{p{2.3cm}p{4.0cm}p{9.1cm}}
\toprule
component & branch & state \\
\midrule
esm\_tools & \texttt{feat/awiesm3-v3.4-co2} & \texttt{cf0a27943} carried the v3.5 CORE3
defaults, the mixing scheme and the IDEMIX forcing paths for the PICAL line. \\
fesom-2.8 & \texttt{main} @ \texttt{82c67a4a} & Contains \#1060, \#1064 and \#1087. Built coupled,
double precision, \texttt{RECOM\_COUPLED=ON}, XIOS on. Hemispheric GM (\#1112) still to be added. \\
oifs-48r1 & \texttt{awiesm3-develop} @ \texttt{f6d1860} (fork) & Five commits past develop's
\texttt{872c4d2}: 2\,m temperature to the ice sheet, and the surface-pressure mass fixer. \\
LPJ-GUESS & \texttt{lpj\_guess\_awiesm3} @ \texttt{179fb0d} & The AWI land-grid line with the
canopy work. \\
OASIS, rnfmap, XIOS & \texttt{local\_combined\_fixes} @ \texttt{898c8b2},
\texttt{enthalpy\_on\_ocean\_side} @ \texttt{a635cc6}, \texttt{main} @ \texttt{9c33506d} & XIOS
re-cloned. \\
\bottomrule
\end{tabular}
\end{center}

Merged upstream during the campaign: FESOM PR \#1054 (the single-precision ocean heat
leak), \#1064 (\texttt{cavity\_gamma\_scale}, \texttt{f3c3d3fa}), \#1060 (ice skin from the
atmosphere, \texttt{57c433e5}), \#1087 (the two CVMix bugs, \texttt{18926b72}), and three CI fixes
of our own that unblocked stacked pull requests in that repository (\#1088, \#1089, \#1091).
\dead{Still open and needed: \#1060 and \#1087.} Both are merged. Open: FESOM \#1103 (IDEMIX
smoothing per unit time), \#1112 (hemispheric GM), and \texttt{release\_evaluation\_tool2} \#7 (the
Gregory snow-enthalpy fix and seven other albedo fixes).

\textbf{On albedo} the 1800\,s production layout is 768 FESOM ranks at 96 per node on 25 nodes; the
1200\,s step needs 33. As of 2026--09--26 the campaign was \emph{moving} to albedo for queue time: a 47-node leg on levante had
waited between 30 seconds and 10 hours to start. The branch point is 2100, the only year
with a complete consistent restart set. Three one-liners kill a first run there, all in
\texttt{HANDOFF\_albedo\_spinup\_continuation.md}: \texttt{h0min} in the runscript aborts
FESOM at the namelist read on any branch except the abandoned
\texttt{feat/leadclose-h0min}; both IDEMIX forcing files must exist under
\texttt{\$\{fesom.pool\_dir\}/forcing/idemix/}; and \texttt{ltos} must go in
\texttt{ecearth.ins.j2}, the file the run actually reads. The standing gate on any IDEMIX
leg is \texttt{Etot\_srf} $0.3211645$ and \texttt{Etot\_bot}
$0.9135067\to0.9179869$\,TW within the first minute of the compute log; the failure
signature is $6.05\times10^{-4}$ with $-1.54\times10^{16}$, which NaN'd the ocean in 100
seconds on 2026--09--24. Two later notes: round 45 found that the \texttt{/ice\_therm/} read does
not check \texttt{iostat}, so an unknown name there ends the read silently rather than aborting
(the abort above was not re-examined against this); and the pool IDEMIX forcing has been unreadable
since 2026--10--01, so the runscripts point \texttt{idemix\_botforc\_file} at a project-local copy.
The spin-up plots of this summary come from \texttt{release\_evaluation\_tool2} branch
\texttt{feat/critical-path-spinup} (per-year cache, annotations, AMOC series).

\section{The path forward}\label{sec:forward}

The plan for the last round, as agreed on 2026--10--03.

\textbf{0. Added the same day: TKE without IDEMIX.} \textbf{PICAL\_crunveg\_gmhemi\_tke}, 2140--2159
from gmhemi1800, one variable changed, runs first (\S\ref{sec:mixing}). Score it on heat by depth,
both hemispheres' ice volume and net TOA before choosing the mixing scheme of the release.

\textbf{1. Hemispheric GM with NH \texttt{K\_GM\_max} 1500.} gmhemi1800 is right in everything but
the Labrador Sea, which it overshoots (ice 0.10--0.13 against 0.16). A higher NH value is the
direct correction, and the GM 1500 arm at 1200\,s already showed a Labrador Sea that convects
without capping.

\textbf{2. Continue the current tuning on the \texttt{-cc} configuration.} FESOM and LPJ-GUESS warm
start, REcoM on in concentration-driven PI mode, OpenIFS cold start, on the fresh install of
\S\ref{sec:stack}. Hemispheric GM must be carried into that build. \emph{Done for one leg:}
\textbf{PICAL\_cc} ran 2140--2144 and is paused at a complete restart for performance work; it costs
6.7 times the physics-only line.

\textbf{3. Document the structural biases of v3.5.} The Southern Ocean cloud deficit, the diffuse
Antarctic pack and its over-responsive drift, the warm stratocumulus coasts, and the warm, dry NH
summer land. Each is measured, each resisted the levers aimed at it or has no physical lever left,
and each is better stated plainly than tuned around.

\medskip
Note what is \emph{not} on this list. DMS / marine CCN was reconsidered as a global cooling lever
and declined; \texttt{h0min} is not a lever in the coupled build; the McPhee ice--ocean flux went
the wrong way in both hemispheres; slowing the Antarctic drift would likely cost SH extent. The
atmospheric lever families were, by measurement, nearly exhausted before these rounds.

\section{Assets}\label{sec:assets}

{\small
\begin{longtable}{p{5.2cm}p{10.8cm}}
\toprule
what & where \\
\midrule
\endfirsthead
\toprule
what & where \\
\midrule
\endhead
Full record, 163 pages & \texttt{report/report.tex} \\
Previous summary (2026--08--24) & \texttt{report/summary\_superseded\_20260926.tex}; the
2026--09--26 version of this one is in git history \\
Albedo handoff & \texttt{HANDOFF\_albedo\_spinup\_continuation.md} \\
Annual TOA / surface / OHC store & \texttt{data/coupled\_annual\_diag.nc}, built by
\texttt{scripts/analysis/coupled\_annual\_store.py} \\
Drift, ice volume, GIOMAS & \texttt{scripts/analysis/ccnice\_regional\_drift.py},
\texttt{ice\_volume\_series.py}, \texttt{giomas\_volume\_targets.py} (takes
\texttt{OBS\_Y0}/\texttt{OBS\_Y1}) \\
Drake Passage, cavity melt & \texttt{scripts/analysis/drake\_passage\_transport.py} and
\texttt{scripts/figures/cavity\_melt\_maps\_mixing.py}; both need the conda env
\texttt{esm-tools\_auto\_tripyview}, and the transport needs
\texttt{PYTHONPATH=/work/ab0246/a270092/software/tripyview} \\
Run and parameter notes & \texttt{notes/RUNS\_AND\_PARAMETERS.md} \\
Labrador history and noise floor & \path{scripts/analysis/labrador_any_mesh.py},
\texttt{labrador\_ice\_vs\_global\_sst.py}, \texttt{labrador\_mld\_noise\_floor.py} \\
GM arms and heat pathways & \path{scripts/analysis/gm_arms_global.py}, \texttt{heat\_pathways.py},
\texttt{gm\_k\_rossby.py}; \texttt{scripts/figures/gm\_k\_maps.py}; tables
\texttt{data/clim/gm\_arms1800\_*.txt} \\
PI bias and warm regions & \path{scripts/analysis/pi_t2m_bias_epoch_corrected.py},
\texttt{pi\_t2m\_bias\_land\_ocean.py}, \texttt{warm\_regions\_diag.py};
\texttt{scripts/figures/pi\_t2m\_bias\_maps.py} \\
Sea-ice budget and drift & \path{scripts/analysis/ice_winter_budget_hemispheres.py} (env
\texttt{SIC}), \texttt{ice\_drift\_vs\_osisaf.py}, \texttt{ice\_wind\_response.py},
\texttt{ice\_drag\_effective.py}; obs under \texttt{/albedo/work/user/jstreffi/obs/} \\
Full \texttt{reval} diagnostics & \texttt{/work/ab0246/a270092/software/release\_evaluation\_tool2},
per-window configs under \texttt{configs/} \\
Configuration & esm\_tools \texttt{feat/awiesm3-v3.4-co2} at \texttt{cf0a27943} \\
\bottomrule
\end{longtable}}

\end{document}
